% Pista ana-23j2-q6 · Anàlisi
% Procedència: PAU Catalunya 2023 · Convocatòria de juny · Sèrie 5 · Qüestió 6
\documentclass[11pt,a4paper]{article}
\usepackage{pista}
\pistainfo{Catalunya 2023}{Convocatòria de juny}{Sèrie 5}

\begin{document}

\section*{\textcolor{blauPista}{Qüestió 6}}

\noindent Sigui la funció $f(x) = \dfrac{ax^{2} + x + b}{x^{2} + 1}$.

\subsection*{\textit{a)} \normalfont Calculeu els valors dels paràmetres $a$ i $b$ si sabem que la gràfica de la funció $f$ té un extrem relatiu en $x = -1$ i passa pel punt $P = \bigl(-2, \frac{13}{5}\bigr)$. \hfill\textnormal{[1,25~punts]}}

\pista{Cada dada de l'enunciat és una equació. «Extrem relatiu en $x = -1$» implica \mbox{$f'(-1) = 0$}, així que la primera cosa que has de fer és calcular la derivada d'un quocient. Quan hagis de resoldre l'equació $f'(-1) = 0$, tingues en compte que només cal igualar a zero el numerador. «Passa per $P$» vol dir $f(-2) = \frac{13}{5}$. Tens, doncs, un sistema de dues equacions amb incògnites $a$ i $b$; la primera et donarà una relació molt neta entre tots dos.}

\subsection*{\textit{b)} \normalfont Per al cas $a = b$, calculeu i classifiqueu els extrems relatius de la funció. \hfill\textnormal{[1,25~punts]}}

\pista{Substitueix la relació $a = b$ a la derivada i simplifica: el numerador es redueix molt. Iguala'l a zero per a trobar els candidats a extrem, i classifica'ls estudiant el signe de $f'$ a banda i banda de cadascun (taula de signes), per a dir on hi ha mínim i on màxim.}

\end{document}
